
\section{Szereg Fouriera: współczynniki, ortogonalność i zbieżność}
\label{sec:dt-series-fourier}
\providecommand{\atlasref}[2]{#2}

Por. \atlasref{PDE-06::thm:pde6-fejer}{RRC-06: średnie Fejéra i szeregi Fouriera}.


Przez $\mathbb T=\R/(2\pi\mathbb Z)$ rozumiemy okrąg z
współrzędną $x\in[-\pi,\pi]$ i utożsamionymi końcami.
Funkcję z $L^1(-\pi,\pi)$ utożsamiamy z jej
$2\pi$-okresowym przedłużeniem, określonym z dokładnością
do zbioru miary zero; wartości na końcach przedziału
nie wpływają na całki. Dla takiej funkcji definiujemy
$\widehat f(k)=(2\pi)^{-1}\int_{-\pi}^{\pi}f(x)e^{-ikx}\,\mathrm d\lambda_1(x)$,
$k\in\mathbb Z$, oraz sumę częściową
$S_Nf(x)=\sum_{|k|\le N}\widehat f(k)e^{ikx}$.
W tym rozdziale $L^2(\mathbb T)$ oznacza zespoloną przestrzeń Hilberta
z iloczynem skalarnym
$\langle f,g\rangle=(2\pi)^{-1}\int_{-\pi}^{\pi}f(x)\overline{g(x)}\,\mathrm d\lambda_1(x)$.
Pojęcie całki Lebesgue'a i podstawowe własności $L^p$
traktujemy jako znane. Na okręgu używamy miary Haara
unormowanej do masy $1$, czyli $\mathrm d\lambda_1/(2\pi)$.
Na torusie $\mathbb T^d$ jej odpowiednikiem jest
$\mathrm d\lambda_d/(2\pi)^d$ na kostce fundamentalnej;
współczynniki Fouriera zawierają zatem czynnik $(2\pi)^{-d}$.
Norma $L^2$ względem nieznormalizowanej miary jest
$(2\pi)^{d/2}$ razy większa. Nie należy utożsamiać tej
normalizacji okresowej z unitarną normalizacją całki
Fouriera na $\R^d$ z definicji \ref{def:fourier-schwartz17}.
Ponieważ miara unormowana na
$\mathbb T$ ma masę $1$, nierówność Cauchy'ego--Schwarza
daje $(2\pi)^{-1}\|f\|_{L^1(-\pi,\pi)}\leq\|f\|_2$,
więc współczynniki Fouriera są określone dla każdego
$f\in L^2(\mathbb T)$.

\begin{lemat}[Ortogonalność i nierówność Bessela]
\label{lem:dt-series-bessel}
Niech $f\in L^2(\mathbb T)$ oraz $N\geq0$ będzie całkowite. Wówczas
$\langle e^{ikx},e^{i\ell x}\rangle=\delta_{k\ell}$,
$\|S_Nf\|_2^2=\sum_{|k|\le N}|\widehat f(k)|^2$ i
$\sum_{|k|\le N}|\widehat f(k)|^2\le\|f\|_2^2$.
\end{lemat}
\begin{proof}
Dla $k=\ell$ całka z jedynki względem miary
unormowanej wynosi $1$. Jeśli $k\ne\ell$, to
\[
 \frac1{2\pi}\int_{-\pi}^{\pi}e^{i(k-\ell)x}\,
 \mathrm d\lambda_1(x)
 =\frac{e^{i(k-\ell)\pi}-e^{-i(k-\ell)\pi}}
 {2\pi i(k-\ell)}=0,
\]
ponieważ $k-\ell$ jest całkowite. Otrzymujemy
ortonormalność. Rozwijając kwadrat normy skończonej
sumy $S_Nf$ i korzystając z tej ortonormalności,
dostajemy $\|S_Nf\|_2^2=\sum_{|k|\le N}|\widehat f(k)|^2$.
Ponadto dla $|\ell|\leq N$
\[
 \langle f-S_Nf,e^{i\ell x}\rangle
 =\widehat f(\ell)-\sum_{|k|\leq N}
 \widehat f(k)\delta_{k\ell}=0.
\]
Zatem $S_N$ jest rzutem ortogonalnym na przestrzeń
rozpiętą przez te wykładniki. W szczególności
$\langle f-S_Nf,S_Nf\rangle=0$, skąd
\[
 \|f\|_2^2=\|S_Nf\|_2^2+\|f-S_Nf\|_2^2.
\]
Nieujemność ostatniego składnika daje nierówność Bessela.
\end{proof}

Nie można jednak wywnioskować z ortogonalności, że $S_Nf\to f$.
Do tego potrzebujemy zupełności układu wykładniczego.

